\subsection{Audit Landing Points: Reproducible Interfaces for Unification Propositions}\label{subsec:group-unification-audit-pointers}

This section introduces no new proof premises but instead systematically grounds the key unification conclusions of this chapter in reproducible outputs and falsifiable checks within the repository,
so that readers can locate the corresponding scripts, generated snippets, and audit certificates without recourse to the interpretive layer.

\paragraph{Verifiable arithmetic facts: Fibonacci--Lie resonance ladder (reproducible exhaustion)}
The stable type counting law \(|X_m|=F_{m+2}\) (in the framework of this chapter) exhibits exact resonances with compact Lie algebra dimensions at small windows (Subsection \ref{subsec:fib-lie-resonance}).
The key part is writing the equation as a Diophantine equation and performing finite-window exhaustion, thereby obtaining a reproducible solution set (independent of interpretive-layer narrative):
\begin{itemize}
  \item \textbf{Generation script:} \texttt{\detokenize{scripts/exp_fibonacci_lie_resonance_ladder.py}}.
  \item \textbf{LaTeX generated outputs:} \texttt{\detokenize{sections/generated/tab_fibonacci_lie_resonance_ladder.tex}},
  \texttt{\detokenize{sections/generated/eq_fib_lie_resonance_solutions.tex}}.
  \item \textbf{Audit JSON:} \texttt{\detokenize{artifacts/export/fibonacci_lie_resonance_ladder.json}}.
\end{itemize}
Any modification to resonance points/solution sets must be reproducible via the above script within the declared scanning window.

\paragraph{Verifiable interface: time factor of the Zeckendorf truncation folding (unique branch/reset events)}
Terminal propositions \ref{thm:terminal-succ-unique-branch-merge}--\ref{thm:terminal-reset-events-sturmian} ground the factor dynamics induced by the time advance \(N\mapsto N+1\) at finite window \(m\) as a finite certificate chain: the unique branch point \(b_m\), the unique merge point \(0^m\), and the reset event set \(\mathcal R_m\) with its binary return time spectrum and counting deviation bounds.
The corresponding verifiable outputs are generated by a single script:
\begin{itemize}
  \item \textbf{Generation script:} \texttt{\detokenize{scripts/exp_fold_zeck_time_factor_succ_reset_audit.py}}.
  \item \textbf{LaTeX generated outputs:} \texttt{\detokenize{sections/generated/tab_fold_zeck_succ_unique_branch_audit_m2_m10.tex}},
  \texttt{\detokenize{sections/generated/tab_fold_zeck_reset_event_audit_m2_m10.tex}},
  \texttt{\detokenize{sections/generated/tab_fold_zeck_reset_event_gaps_m6.tex}}.
  \item \textbf{Audit JSON:} \texttt{\detokenize{artifacts/export/fold_zeck_time_factor_succ_reset_audit.json}}.
\end{itemize}

\paragraph{Verifiable interface: information-theoretic constants of \texorpdfstring{$\mathrm{Fold}^{\mathrm{bin}}_6$}{Foldbin6} (information loss / entropy scale)}
Under the uniform microstate baseline \(U\sim\mathrm{Unif}(\{0,\dots,63\})\), let \(W=\mathrm{Fold}^{\mathrm{bin}}_6(U)\).
Terminal proposition \ref{thm:terminal-foldbin6-information-loss-constants} expresses the single-step ``information loss'' \(I_{\mathrm{lost}}=\mathbb{E}[\log_2 d^{\mathrm{bin}}_6(W)]\), the output entropy \(H_2(W)\), and the KL gap from the uniform distribution as closed-form constants depending only on the fiber degeneracy spectrum, with certificate output from finite enumeration:
\begin{itemize}
  \item \textbf{Generation script:} \texttt{\detokenize{scripts/exp_foldbin6_information_loss_constants.py}}.
  \item \textbf{LaTeX generated output:} \texttt{\detokenize{sections/generated/eq_foldbin6_information_loss_constants.tex}}.
  \item \textbf{Audit JSON:} \texttt{\detokenize{artifacts/export/foldbin6_information_loss_constants.json}}.
\end{itemize}

\paragraph{Verifiable interface: affine geometry classification of fibers from the \texorpdfstring{$\mathbb{F}_2^6$}{F2^6} perspective}
Terminal proposition \ref{thm:terminal-foldbin6-fiber-affine-geometry} establishes the affine geometry type of \(F_w=F_6^{-1}(w)\subset\mathbb{F}_2^6\) as an independent auditable rigidity interface: at window-\(6\), only \(11/21\) fibers are affine flat (of which \(8\) are \(1\)-flat and \(3\) are \(2\)-flat), while the remaining fibers are nonlinear subsets (and accordingly negate any ``linear kernel'' simplification, Corollary \ref{cor:terminal-foldbin6-no-linear-kernel}).
The corresponding per-fiber certificate output is generated by a single script:
\begin{itemize}
  \item \textbf{Generation script:} \texttt{\detokenize{scripts/exp_foldbin6_fiber_affine_geometry.py}}.
  \item \textbf{LaTeX generated output:} \texttt{\detokenize{sections/generated/tab_foldbin6_fiber_affine_geometry.tex}}.
  \item \textbf{Audit JSON:} \texttt{\detokenize{artifacts/export/foldbin6_fiber_affine_geometry.json}}.
\end{itemize}

\paragraph{Verifiable interface: the \(18\oplus 3\) splitting at \(m=6\) and the \(B_3/C_3\) unification seed}
Under the boundary word tower framework, the endpoint fiber at \(m=6\) gives three boundary words (see \texttt{\detokenize{sections/generated/eq_boundary_tower_sm12.tex}}), and the \(18\oplus 3\) splitting \(X_6=X_6^{\mathrm{cyc}}\sqcup X_6^{\mathrm{bdry}}\) restricts the \(21\)-dimensional unification seed to \(B_3/C_3\) (Subsection \ref{subsec:bdry-tower-zeck-gut}).
To avoid manual pairing, this paper further grounds the seed as root dictionary and Weyl group orbit size verification outputs:
\begin{itemize}
  \item \textbf{Generation script:} \texttt{\detokenize{scripts/exp_fold6_b3c3_seed.py}}.
  \item \textbf{LaTeX generated outputs:} \texttt{\detokenize{sections/generated/eq_fold6_b3c3_seed.tex}},
  \texttt{\detokenize{sections/generated/tab_fold6_b3c3_root_dictionary_B3.tex}},
  \texttt{\detokenize{sections/generated/tab_fold6_b3c3_root_dictionary_C3.tex}}.
  \item \textbf{Audit JSON:} \texttt{\detokenize{artifacts/export/fold6_b3c3_seed.json}} (containing the \(18\oplus 3\), \(6+12\) fingerprint, both dictionaries, and the \(|W|=48\) orbit verification histogram).
\end{itemize}

\begin{corollary}[Rigid selection of \texorpdfstring{$\mathfrak g_{\mathrm{SM}}$}{gSM} by the minimal resonance triple of the boundary tower]\label{cor:bdry-tower-minimal-resonance-triple-sm}
On the boundary word three-layer tower \(m=4,6,8\), finite exhaustive output gives
\[
\bigl(|X^{\mathrm{bdry}}_4|,\ |X^{\mathrm{bdry}}_6|,\ |X^{\mathrm{bdry}}_8|\bigr)=(1,3,8)
\]
(see \texttt{\detokenize{sections/generated/eq_boundary_tower_sm12.tex}}).
If one requires that the continuous envelope of the boundary modes be a compact connected reductive Lie group, and that its Lie algebra decomposes as a direct sum of three reductive factors with dimensions corresponding layer by layer to \((1,3,8)\),
then this Lie algebra is unique up to isomorphism, and must be
\[
\mathfrak g_{\mathrm{SM}}\ \cong\ \mathfrak u(1)\ \oplus\ \mathfrak{su}(2)\ \oplus\ \mathfrak{su}(3).
\]
\end{corollary}
\begin{proof}
Compact reductive Lie algebras decompose as a center plus a semisimple part.
The compact simple factor of dimension \(3\) is unique up to isomorphism as \(\mathfrak{su}(2)\), the compact simple factor of dimension \(8\) is uniquely \(\mathfrak{su}(3)\), and the abelian factor of dimension \(1\) can only be \(\mathfrak u(1)\).
Juxtaposing the three block dimensions yields the conclusion. \qed
\end{proof}

\begin{remark}[Non-circular selection of \texorpdfstring{$\mathfrak{su}(3)$}{su3} and \texorpdfstring{$\mathfrak{su}(2)$}{su2} via Coxeter/Weyl data]\label{rem:window6-coxeter-weyl-select-su3-su2}
Under the root dictionary framework of \S\ref{subsec:bdry-tower-zeck-gut}, the Weyl group of window-\(6\) satisfies
\(
W\cong W(B_3)\cong W(C_3)\cong (\ZZ_2)^3\rtimes S_3
\).
The Coxeter classification shows that among rank-\(2\) irreducible reflection groups, the group \(S_3\) uniquely corresponds to \(A_2\)-type Weyl data; the rank-\(1\) reflection group \(\ZZ_2\) uniquely corresponds to \(A_1\)-type Weyl data.
Therefore, once \(S_3\) and \(\ZZ_2\) are intrinsically produced as auditable finite reflection symmetries, juxtaposing with the dimensional interface of Corollary \ref{cor:bdry-tower-minimal-resonance-triple-sm} yields a non-circular localization of \(\mathfrak{su}(3)\) and \(\mathfrak{su}(2)\).
\end{remark}

\begin{remark}[Weyl two-orbit compression and the two degrees of freedom of ``coupling normalization'' (window-\(6\))]\label{rem:fold6-weyl-two-orbit-two-dof}
Corollary \ref{cor:fold6-weyl-two-orbit-compression} shows that any Weyl-invariant observable on the cyclic sector \(X^{\mathrm{cyc}}_6\) is determined by only two orbit constants (short root/long root).
Therefore, if one views the weight function on root directions, or equivalently the \(W\)-invariant quadratic form parametrization, as an abstract model of ``gauge coupling normalization choice,''
then the unification seed at window-\(6\) allows at most two degrees of freedom (short/long classes) at this level.
Furthermore, when this weight function must also be compatible with the boundary sector, or with the common refinement of the orbit--Levi partition,
any ``alignment/unification'' constraint must reduce to a relation between these two degrees of freedom, rather than an external continuous fine-tuning assumption.
\end{remark}

\paragraph{Parallel verifiable interface: the \texorpdfstring{$C_6$}{C6} orbit decomposition at \(m=6\) and the \texorpdfstring{$21=15\oplus 3\oplus 3$}{21=15+3+3} Levi skeleton candidate}
At the same window-6 anchor, the cyclic rotation \(C_6\) acts closedly on the cyclic sector \(X_6^{\mathrm{cyc}}\) and induces a rigid orbit multiset \(\{1,2,3,6,6\}\) (see the paragraph on \texorpdfstring{$C_6$}{C6} orbit rigidity in Subsection \ref{subsec:bdry-tower-zeck-gut}).
This rigid orbit decomposition gives a canonical \(15\oplus 3\) subdivision, which together with the boundary sector's \(3\) forms the \(21=15\oplus 3\oplus 3\) dimensional skeleton,
thereby placing a Pati--Salam-type minimal semisimple Levi envelope as a parallel candidate grounded in auditable output.
Additionally, the dyadic folding \(\mathrm{Fold}^{\mathrm{bin}}_6\) forces a \(\delta=34\) double-layer covering on the boundary sector and induces a verifiable \(\ZZ_2\) sheet parity (see the boundary sector double-layer covering construction in Subsection \ref{subsec:bdry-tower-zeck-gut}).
\begin{itemize}
  \item \textbf{Generation script:} \texttt{\detokenize{scripts/exp_window6_c6_orbit_patisalam_seed.py}}.
  \item \textbf{LaTeX generated outputs:} \texttt{\detokenize{sections/generated/eq_window6_c6_orbit_decomposition.tex}},
  \texttt{\detokenize{sections/generated/eq_window6_c6_character_decomposition.tex}},
  \texttt{\detokenize{sections/generated/tab_window6_c6_orbit_decomposition.tex}},
  \texttt{\detokenize{sections/generated/tab_window6_patisalam_9633_partition.tex}},
  \texttt{\detokenize{sections/generated/eq_fold6_bin_uplift_delta_set.tex}},
  \texttt{\detokenize{sections/generated/tab_fold6_bin_uplift_choice_collapse.tex}},
  \texttt{\detokenize{sections/generated/tab_fold6_boundary_sheet_pairs.tex}},
  \texttt{\detokenize{sections/generated/tab_foldbin_boundary_lift_m6_m8.tex}}.
  \item \textbf{Audit JSON:} \texttt{\detokenize{artifacts/export/window6_c6_orbit_patisalam_seed.json}} (containing \(C_6\) orbits, the \(9+6+3+3\) four-block decomposition, \(\Delta_6=\{0,21,34,55\}\), the uplift-choice collapse partition \(\Xi(w)\in\{2,3,4\}\), \(w_6=1\Rightarrow\{0,34\}\), explicit preimage pairs of the three boundary words, and the boundary-lift three-layer/two-layer stability scan outputs for \(m=6,7,8\)).
\end{itemize}

\begin{remark}[Cross-resolution extension of sheet parity is not automatic: additional protocol gates are needed to re-lock the double-layer covering]\label{rem:window6-sheet-parity-nonpersistent}
By Table \ref{tab:foldbin_boundary_lift_m6_m8}, the boundary sector at \(m=6\) under dyadic uplift exhibits a two-layer covering pattern \(\{0,34\}\);
but under the same framework, when the resolution increases to \(m=7,8\), the uplift deltas become \(\{0,55,89\}\) and \(\{0,89,144\}\) respectively, with the number of covering layers turning to three.
Therefore the \(\ZZ_2\) sheet parity of window-\(6\) is a hard structural fact specific to the minimal anchor; to promote it as a cross-scale residual selection law requires imposing additional auditable protocol gates to re-lock the double-layer covering (e.g., further constraints on uplift-choice or the introduction of additional readout registers).
\end{remark}

\begin{remark}[Falsifiable boundary under local uplift framework: \texorpdfstring{$\delta\neq 34$}{delta!=34} necessarily forces non-local externalization]\label{rem:window6-local-uplift-falsifiable-boundary}
Corollary \ref{cor:terminal-window6-local-uplift-admissibility} shows that under the strong local framework where ``uplift is realized by label-preserving geometric automorphisms,'' the only admissible displacement besides the trivial one is \(\delta=34\).
Therefore, any candidate extension requiring \(\delta=21\) or \(\delta=55\) type tail branch triggering, if still claiming to operate under this local geometric uplift framework, contradicts the geometric stabilizer rigidity;
to realize it, one must choose between: abandoning the locality of ``label-preserving and realized by \(\mathrm{Aut}(Q_6)\),'' or introducing additional projection gates/protocol constraints to alter the realizable set of uplift.
\end{remark}

\begin{remark}[``Double-seed bifurcation'' of dimension \(21\) and falsifiable discriminants]\label{rem:window6-dim21-double-seed-bifurcation}
Window-\(6\) simultaneously outputs at the same \(m=6\) anchor two mutually independent \(21\)-dimensional continuous envelope candidates, both verifiable by finite exhaustive certificates:
one locks the simple type seed to \(B_3/C_3\) via the \(18\oplus 3\) root--Cartan dictionary;
the other locks the Levi-type skeleton to \(\mathfrak{su}(4)\oplus\mathfrak{su}(2)_L\oplus\mathfrak{su}(2)_R\) via \(C_6\) orbit rigidity.
Therefore, ``the choice of unification candidate'' in the framework of this paper is not an interpretive prior but a falsifiable discrimination task: additional, equally auditable discriminants must be introduced (such as twisted spectra/anomaly signatures and block alignment constraints at higher resolution, see \S\ref{subsec:group-unification-spectral-alignment} and its script/certificate chain) to distinguish the two types of seeds.
\end{remark}

\begin{corollary}[Levi rigidity locking: the Pati--Salam-type skeleton forced by \texorpdfstring{$21=15\oplus 3\oplus 3$}{21=15+3+3}]\label{cor:window6-levi-rigidity-pati-salam}
Under the orbit--Levi framework of \S\ref{subsec:bdry-tower-zeck-gut}, if one further imposes the minimal continuous assumption: the continuous envelope of window-\(6\) admits only compact semisimple Lie algebras,
and the \(21=15\oplus 3\oplus 3\) dimensional skeleton corresponds to a direct sum decomposition of three compact simple ideals,
then this Levi skeleton is unique up to isomorphism and must be isomorphic to
\[
\mathfrak g^{\mathrm{Levi}}_6\ \cong\ \mathfrak{su}(4)\ \oplus\ \mathfrak{su}(2)_L\ \oplus\ \mathfrak{su}(2)_R.
\]
\end{corollary}
\begin{proof}
The compact simple Lie algebra of dimension \(3\) is unique up to isomorphism, namely \(\mathfrak{su}(2)\cong\mathfrak{so}(3)\).
On the other hand, compact simple Lie algebra dimensions can only take the classical type values
\(\dim\mathfrak{su}(n)=n^2-1\), \(\dim\mathfrak{so}(n)=\frac{n(n-1)}2\), \(\dim\mathfrak{sp}(n)=n(2n+1)\)
and the exceptional type values \(14,52,78,133,248\), etc.
Direct checking shows: the unique compact simple isomorphism class with \(\dim=15\) is \(\mathfrak{su}(4)\cong\mathfrak{so}(6)\).
Juxtaposing the three block dimensions yields the conclusion. \qed
\end{proof}

\begin{remark}[\texorpdfstring{$15=1\oplus 2\oplus 6\oplus 6$}{15=1+2+6+6} orbit dictionary and the non-external alignment of \texorpdfstring{$\mathfrak{su}(4)\supset\mathfrak{su}(3)\oplus\mathfrak u(1)$}{su4 contains su3+u1}]\label{rem:window6-su4-su3u1-orbit-dictionary}
Under the setting of Corollary \ref{cor:window6-levi-rigidity-pati-salam}, the \(15\)-dimensional block can be identified up to isomorphism as the adjoint component of \(\mathfrak{su}(4)\).
Its standard maximal subalgebra \(\mathfrak{su}(3)\oplus\mathfrak u(1)\subset\mathfrak{su}(4)\) induces the adjoint decomposition
\[
\mathfrak{su}(4)\ \cong\ (\mathfrak{su}(3)\oplus\mathfrak u(1))\ \oplus\ (\mathbf3\oplus\bar{\mathbf3}),
\]
and further decomposes \(\mathfrak{su}(3)\) into Cartan and step directions:
\(
\mathfrak{su}(3)=\mathfrak h\oplus\mathfrak n,\ \dim\mathfrak h=2,\ \dim\mathfrak n=6
\).
Therefore at the pure dimensional level one obtains
\[
15\ =\ 1\ \oplus\ 2\ \oplus\ 6\ \oplus\ 6,
\]
consistent with the orbit size multiset of the \(15\)-element cyclic block read off from the \(C_6\) orbit rigidity at window-\(6\) (see the preceding paragraph and the generated Table \ref{tab:window6_c6_orbit_decomposition}).
In the framework of this paper, this alignment means: once ``orbit size dictionary'' and ``compact semisimple Levi skeleton'' compatibility is required, the appearance of \(\mathfrak{su}(4)\supset\mathfrak{su}(3)\oplus\mathfrak u(1)\) needs no additional embedding choice.
\end{remark}

\begin{corollary}[Embedding rigidity of \texorpdfstring{$\mathfrak{su}(4)\supset \mathfrak{su}(3)\oplus\mathfrak u(1)$}{su4 contains su3+u1} (conjugacy uniqueness)]\label{cor:window6-su4-embedding-rigidity}
Under the setting of Corollary \ref{cor:window6-levi-rigidity-pati-salam}, identify the \(15\)-dimensional block as the adjoint component of \(\mathfrak{su}(4)\).
Then the maximal-rank subalgebra chain compatible with the orbit dictionary \(15=1\oplus 2\oplus 6\oplus 6\),
\[
\mathfrak{su}(3)\oplus\mathfrak u(1)\ \subset\ \mathfrak{su}(4),
\]
is conjugacy-unique under the action of \(\mathrm{Aut}(\mathfrak{su}(4))\).
Therefore the \(\mathfrak u(1)\) direction as the commutant direction of \(\mathfrak{su}(3)\) (i.e., their relative position within \(\mathfrak{su}(4)\)) is rigid data under this discrete--continuous alignment framework, not a freely interchangeable external choice.
\end{corollary}
\begin{proof}
In \(\mathfrak{su}(4)\), the standard embedding of \(\mathfrak{su}(3)\oplus\mathfrak u(1)\) can be taken as the diagonal block stabilizer (equivalently: the Lie algebra of the stabilizer of a complex line in \(\mathbb{C}^4\) under the \(SU(4)\) action).
Since \(SU(4)\) acts transitively on the complex projective space \(\mathbb{CP}^3\), all such subalgebras are conjugate under inner automorphisms; hence conjugacy-unique under \(\mathrm{Aut}(\mathfrak{su}(4))\). \qed
\end{proof}

\begin{corollary}[Levi rigidity of 6D window-\(6\) forces a Standard Model-type residual algebra]\label{cor:window6-levi-rigidity-sm-residual}
On top of the Levi skeleton of Corollary \ref{cor:window6-levi-rigidity-pati-salam}
\(
\mathfrak g^{\mathrm{Levi}}_6\cong \mathfrak{su}(4)\oplus\mathfrak{su}(2)_L\oplus\mathfrak{su}(2)_R
\),
juxtapose the embedding rigidity of \(\mathfrak{su}(4)\supset \mathfrak{su}(3)\oplus\mathfrak u(1)\) (Corollary \ref{cor:window6-su4-embedding-rigidity})
with the geometric \(\ZZ_2\) of window-\(6\) (Corollary \ref{cor:terminal-foldbin6-geo-stabilizer-z2}, see also Corollary \ref{cor:terminal-window6-local-uplift-admissibility} for the \(\delta=34\) selection law).
Then under the minimal framework introducing no additional continuous degrees of freedom, there exists a one-dimensional fixed-point direction \(\mathfrak u(1)_{\mathrm{eff}}\subset \mathfrak{su}(2)_R\) such that
\[
\mathfrak{su}(3)\ \oplus\ \mathfrak{su}(2)_L\ \oplus\ \mathfrak u(1)_{\mathrm{eff}}
\subseteq \mathfrak g^{\mathrm{Levi}}_6.
\]
Thus the Standard Model-type non-abelian skeleton \(\mathfrak{su}(3)\oplus\mathfrak{su}(2)_L\) and an effective abelian direction can be intrinsically selected by the combinatorial/representation-theoretic rigidity of window-\(6\), without relying on the circular premise of ``specifying the Standard Model a priori.''
\end{corollary}
\begin{proof}
By Corollary \ref{cor:window6-su4-embedding-rigidity}, the maximal-rank embedding \(\mathfrak{su}(3)\oplus\mathfrak u(1)\) compatible with the orbit dictionary within \(\mathfrak{su}(4)\) is conjugacy-unique;
therefore within \(\mathfrak g^{\mathrm{Levi}}_6\) there exists a subalgebra locked by discrete outputs:
\(
\mathfrak{su}(3)\oplus\mathfrak u(1)\oplus\mathfrak{su}(2)_L\oplus\mathfrak{su}(2)_R
\).
On the other hand, Corollary \ref{cor:terminal-foldbin6-geo-stabilizer-z2} gives the unique nontrivial geometric involution at window-\(6\) (equivalently: an auditable \(\ZZ_2\) sheet parity constraint);
viewing it under the minimal continuous envelope framework as a nontrivial involutive automorphism of the right-handed block \(\mathfrak{su}(2)_R\), its fixed-point Lie algebra must be one-dimensional.
Indeed \(\mathfrak{su}(2)\cong\mathfrak{so}(3)\), and any nontrivial involution on \(\mathfrak{so}(3)\) is conjugate to a reflection about an axis, whose fixed point is precisely the rotation direction about that axis, hence isomorphic to \(\mathfrak u(1)\).
Taking this fixed-point direction as \(\mathfrak u(1)_{\mathrm{eff}}\subset \mathfrak{su}(2)_R\) yields the conclusion. \qed
\end{proof}

\paragraph{Verifiable interface: the edge--flux skeleton of window-6 (6D hypercube local perturbation \(\times\) \(9+6+3+3\) four-block decomposition)}
Terminal proposition \ref{thm:terminal-window6-edge-flux-skeleton} pushes the one-step local perturbation of the hypercube graph \(Q_6\) onto the four-block Levi skeleton of \(X_6\),
obtaining a fully explicit and verifiable four-block edge-count matrix \(E\), and from it constructs the coarse-grained four-state Markov kernel \(T\) and its spectral fingerprint (the \((t-1)\) factor decomposition of the characteristic polynomial).
The same certificate also gives the cut size and steady-state flux constant \(\Phi_R=\frac1{12}\) for the boundary sector \(U_R\), and outputs the stationary distribution/detailed balance certificate for the pushforward Markov kernel \(P_6\) on \(X_6\) (Theorem \ref{thm:terminal-foldbin6-pushforward-markov}):
\begin{itemize}
  \item \textbf{Generation script:} \texttt{\detokenize{scripts/exp_window6_edge_flux_skeleton.py}}.
  \item \textbf{LaTeX generated outputs:} \texttt{\detokenize{sections/generated/eq_window6_pushforward_markov_kernel.tex}},
  \texttt{\detokenize{sections/generated/eq_window6_edge_flux_skeleton.tex}}.
  \item \textbf{Spectral gap audit (\(\lambda_2\) and \(\lambda_\star\)):}
  \texttt{\detokenize{scripts/exp_window6_pushforward_markov_spectral_gap.py}} generates the snippet
  \texttt{\detokenize{sections/generated/eq_window6_pushforward_markov_spectral_gap.tex}}
  and the certificate \texttt{\detokenize{artifacts/export/window6_pushforward_markov_spectral_gap.json}}.
  \item \textbf{Audit JSON:} \texttt{\detokenize{artifacts/export/window6_pushforward_markov_kernel.json}},
  \texttt{\detokenize{artifacts/export/window6_edge_flux_skeleton.json}}.
\end{itemize}

\begin{remark}[The edge--flux skeleton as an iterable coarse-graining operator candidate]\label{rem:window6-edge-flux-coarse-graining-iterable}
Theorem \ref{thm:terminal-window6-edge-flux-skeleton} gives at \(m=6\) an explicit pushforward/coarse-graining from microscopic edge perturbations to the \(9\oplus 6\oplus 3\oplus 3\) Levi partition.
In the multi-resolution framework, this can be understood as a finite functor prototype ``from microscale to block level'': its input is an auditable edge count, and its output is the block-level kernel \(T\) and its spectral fingerprint;
the portability and iterability of this construction (cross-\(m\) renormalization structure) can be further tested and falsified by maintaining block compatibility and spectral fingerprint stability at higher \(m\).
\end{remark}

\begin{remark}[Normalization-invariant status of the boundary flux constant \texorpdfstring{$\Phi_R=\tfrac1{12}$}{PhiR=1/12}]\label{rem:window6-boundary-flux-phiR-normalization}
The steady-state flux constant \(\Phi_R=\frac1{12}\) given by Theorem \ref{thm:terminal-window6-edge-flux-skeleton} is entirely determined by the finite edge-count certificate,
and under labeling rigidity (Corollary \ref{cor:terminal-gamma6-multiplicity-rigidity}) admits no relabeling freedom.
Therefore \(\Phi_R\) can serve as a normalization-invariant of the window-\(6\) coarse-grained dynamics: any realization claiming to be isomorphic to this edge--flux interface that fails to reproduce \(\Phi_R\) can be directly falsified by finite certificate.
\end{remark}

\begin{corollary}[Compactness principle derived from the reversible detailed balance kernel (window-\(6\))]\label{cor:window6-p6-compactness-principle}
In the notation of Theorem \ref{thm:terminal-foldbin6-pushforward-markov}, \(P_6\) satisfies detailed balance with respect to the stationary distribution \(\pi_6\), and hence the induced Markov operator
\[
(\mathsf{T}_6 f)(w):=\sum_{w'\in X_6} P_6(w,w')\,f(w')
\]
is self-adjoint on the Hilbert space \(\ell^2(X_6,\pi_6)\).
Let its commutant algebra be
\[
\mathcal{C}_6:=\bigl\{U\in \mathrm{End}(\ell^2(X_6,\pi_6)):\ U\mathsf{T}_6=\mathsf{T}_6U\bigr\}.
\]
Then \(\mathcal{C}_6\) is a finite-dimensional \(*\)-algebra, and its unitary group
\(
\mathrm{U}(\mathcal{C}_6):=\{U\in\mathcal{C}_6:\ U^*U=\mathrm{id}\}
\)
is a compact Lie group.
In particular, any continuous gauge symmetry defined by ``preserving the reversible readout kernel \(P_6\)'' must be a compact group.
\end{corollary}
\begin{proof}
Detailed balance is equivalent to \(\mathsf{T}_6\) being self-adjoint with respect to the inner product \(\langle f,g\rangle_{\pi_6}:=\sum_{w\in X_6} f(w)\overline{g(w)}\,\pi_6(w)\).
Since \(|X_6|=21\), the space \(\ell^2(X_6,\pi_6)\) is finite-dimensional, so \(\mathrm{U}(\mathcal{C}_6)\) is a closed subgroup of \(U(21)\), hence compact. \qed
\end{proof}

\paragraph{Verifiable interface: the \texorpdfstring{$\mathrm{Aut}(Q_6)$}{Aut(Q6)} geometric stabilizer (unique \texorpdfstring{$\ZZ_2$}{Z2}) and the \texorpdfstring{$\delta=34$}{delta=34} mask fingerprint}
Beyond the fold-gauge group of ``arbitrary permutations within fibers,'' the main text in Subsection \ref{subsec:bdry-tower-zeck-gut} further introduces a geometric-level stabilizer:
viewing \(F_6(\omega)=\mathrm{Fold}^{\mathrm{bin}}_6(\mathrm{int}_6(\omega))\) as a labeling function on the vertex set of the hypercube graph \(Q_6\), and filtering within \(\mathrm{Aut}(Q_6)\) for elements that commute with it, one obtains the unique surviving nontrivial \(\ZZ_2\) symmetry (Theorem \ref{thm:foldbin6-geo-stabilizer-z2}), which in binary encoding carries the mask integer \(34=100010_2\) and induces a \(\pm 34\) offset on the half-cube (Corollary \ref{cor:foldbin6-geo-mask-34}).
\begin{itemize}
  \item \textbf{Generation script:} \texttt{\detokenize{scripts/exp_foldbin6_geo_stabilizer.py}}.
  \item \textbf{LaTeX generated output:} \texttt{\detokenize{sections/generated/eq_foldbin6_geo_stabilizer.tex}}.
  \item \textbf{Audit JSON:} \texttt{\detokenize{artifacts/export/foldbin6_geo_stabilizer.json}} (containing the exhaustive certificate for \(|\mathrm{Stab}_{\mathrm{geo}}|=2\), the \((\mathrm{mask},\mathrm{perm})\) parametrization of the generator, fixed-point/support counts, and the binary encoding difference histogram).
  \item \textbf{Geometric audit interface (\(\sigma_{\mathrm{geo}}\) orbit-charge splitting within fibers).}
  \texttt{\detokenize{scripts/exp_foldbin6_geo_orbit_charge.py}} generates the table
  \texttt{\detokenize{sections/generated/tab_foldbin6_geo_orbit_charge.tex}}
  and the certificate \texttt{\detokenize{artifacts/export/foldbin6_geo_orbit_charge.json}},
  giving \((f_x,p_x)\) (fixed-point count/two-cycle count) for each \(x\in X_6\) and a seven-class histogram, directly verifiable against Corollary \ref{cor:terminal-foldbin6-geo-orbit-charge}.
  \item \textbf{Dynamical audit interface (explicit counterexample for strong lumpability failure).}
  \path{scripts/exp_foldbin6_strong_lumpability_counterexample.py} generates the snippet
  \texttt{\detokenize{sections/generated/eq_foldbin6_strong_lumpability_counterexample.tex}}
  and the certificate \texttt{\detokenize{artifacts/export/foldbin6_strong_lumpability_counterexample.json}},
  giving the one-step neighborhood stable-type count distribution discrepancy between two same-fiber points, directly verifiable against the assertion ``the Foldbin6 fiber partition does not have strong lumpability.''
  \item \textbf{Geometric audit interface (three-valued law of intra-fiber minimum Hamming distance).}
  \path{scripts/exp_foldbin6_fiber_hamming_min_distance.py} generates the table
  \path{sections/generated/tab_foldbin6_fiber_hamming_min_distance.tex},
  the count histogram snippet \path{sections/generated/eq_foldbin6_fiber_hamming_min_distance_value_counts.tex},
  and the certificate \path{artifacts/export/foldbin6_fiber_hamming_min_distance.json},
  giving \(\delta_{\min}(w)\) and the preimage list for each \(w\in X_6\), directly verifiable against the \(\{2,3,5\}\) three-valued law of Theorem \ref{thm:terminal-foldbin6-fiber-hamming-three-valued}.
  \item \textbf{Geometric audit interface (rigidity of the type adjacency graph \texorpdfstring{$\Gamma_6$}{Gamma6}).}
  \texttt{\detokenize{scripts/exp_window6_type_adjacency_graph_rigidity.py}} generates the snippet
  \texttt{\detokenize{sections/generated/eq_window6_type_adjacency_graph_rigidity.tex}}
  and the certificate \texttt{\detokenize{artifacts/export/window6_type_adjacency_graph_rigidity.json}},
  giving \(|V|,|E|\), the degree histogram, connectivity, diameter, simple spectrum, and \(\mathrm{Aut}(\Gamma_6)=\{1\}\) as verifiable conclusions (Proposition \ref{prop:terminal-gamma6-rigidity} and Corollary \ref{cor:terminal-gamma6-diameter-simple-spectrum}).
\end{itemize}

\begin{corollary}[Finite completeness of 6D observable determination: a decidable isomorphism template from \texorpdfstring{$(\pi_6,P_6,\sigma_{\mathrm{geo}})$}{(pi6,P6,sigma_geo)}]\label{cor:terminal-window6-finite-completeness-template}
Under the window-\(6\) audit framework, let
\[
\mathfrak{A}_6:=(X_6,\pi_6,P_6,\langle\sigma_{\mathrm{geo}}\rangle),
\]
where \(\pi_6\) and \(P_6\) are given by Theorem \ref{thm:terminal-foldbin6-pushforward-markov}, and \(\langle\sigma_{\mathrm{geo}}\rangle\) by Theorem \ref{thm:foldbin6-geo-stabilizer-z2}.
Then for any other realization that outputs isomorphic data at the \(m=6\) level (e.g., from any 6D cut-and-project instance satisfying Proposition \ref{prop:terminal-cut-project-to-fold-hist}),
if it produces the same triple \((\pi_6,P_6,\langle\sigma_{\mathrm{geo}}\rangle)\) under some labeling, then this labeling is uniquely determined by the weighted adjacency rigidity of Corollary \ref{cor:terminal-gamma6-multiplicity-rigidity};
hence ``whether isomorphic'' is a finitely decidable problem, and any deviation can be directly falsified by finite certificate.
\end{corollary}
\begin{proof}
Given \(\pi_6\) and \(P_6\), one can recover the pushforward adjacency count \(A_6\) and the corresponding weighted adjacency multiplicity matrix \(A=(A_{uv})\) (see the defining equation of Theorem \ref{thm:terminal-foldbin6-pushforward-markov}).
Corollary \ref{cor:terminal-gamma6-multiplicity-rigidity} rules out any nontrivial relabeling that preserves \(A\), so the labeling is unique.
The geometric \(\ZZ_2\) stabilizer is uniquely given by Theorem \ref{thm:foldbin6-geo-stabilizer-z2}.
Therefore, comparing two realizations reduces to equality checks on finite objects, and the conclusion holds. \qed
\end{proof}

\begin{remark}[Auditable uniqueness of structurally defined objects (window-\(6\))]\label{rem:terminal-window6-structure-defined-uniqueness}
By Proposition \ref{prop:terminal-gamma6-rigidity} and Corollary \ref{cor:terminal-gamma6-multiplicity-rigidity}, any structure that depends only on the edge pushforward data of \(Q_6\) (equivalently: only on the weighted adjacency multiplicity matrix \(A\), or on \(P_6\), edge--flux skeleton, etc., constructed from it) has no nontrivial relabeling freedom up to isomorphism.
Therefore, any additional objects defined functorially by ``pushforwarding local perturbations of \(Q_6\) to \(X_6\)'' (partitions, coarse-grainings, spectral fingerprints, arithmetic invariants, etc.) can be regarded as label-canonical and auditable finite objects; any claimed alternative realization that is inconsistent on these finite outputs can be directly falsified by finite certificate.
\end{remark}

\begin{remark}[Geometric rigidity eliminates relabeling freedom: canonical localization of the three factors (window-\(6\))]\label{rem:window6-no-relabeling-factor-localization}
Proposition \ref{prop:terminal-gamma6-rigidity} gives \(\mathrm{Aut}(\Gamma_6)=\{1\}\), so any type relabeling that preserves the local adjacency structure must be the identity;
therefore the boundary/cyclic partition, root dictionary, and \(P_6\) defined by the edge pushforward of \(Q_6\) are all canonical in the labeling sense (see also Remark \ref{rem:terminal-window6-structure-defined-uniqueness}).
Juxtaposing with the unique geometric stabilizer \(\langle\sigma_{\mathrm{geo}}\rangle\cong\ZZ_2\) of Theorem \ref{thm:foldbin6-geo-stabilizer-z2} shows that
\(\sigma_{\mathrm{geo}}\) cannot be absorbed by any relabeling.
Hence the \((1,3,8)\) dimensional interface from the boundary word tower and the
\(\mathfrak u(1)\oplus\mathfrak{su}(2)\oplus\mathfrak{su}(3)\) decomposition of Corollary \ref{cor:bdry-tower-minimal-resonance-triple-sm} have canonical localization under this audit framework:
any claimed ``factor interchange/level reallocation'' that preserves the local adjacency and pushforward kernel data must conflict with \(\mathrm{Aut}(\Gamma_6)=\{1\}\).
\par
Furthermore, since the geometric stabilizer is only \(\langle\sigma_{\mathrm{geo}}\rangle\cong\ZZ_2\) (Theorem \ref{thm:foldbin6-geo-stabilizer-z2}),
the geometric side cannot induce any nontrivial relabeling freedom at the type level;
therefore any nontrivial symmetry appearing at the \(m=6\) readout level (such as the boundary sheet-flip and its central embedding, Corollary \ref{cor:terminal-delta-parity-rule} and Corollary \ref{cor:fold-bin-boundary-center-z2})
can only come from the intra-fiber permutations of the fold-gauge group \(\mathcal{G}^{\mathrm{bin}}_6\), not from the geometric automorphisms of \(\mathrm{Aut}(Q_6)\).
\end{remark}

\paragraph{Verifiable interface: exact stability polynomial of \texorpdfstring{$\mathrm{Fold}^{\mathrm{bin}}_6$}{Foldbin6} under independent bit-flip noise}
Terminal proposition Theorem \ref{thm:terminal-foldbin6-bitflip-stability-polynomial} expresses the stability probability of the window-\(6\) stable type readout under an i.i.d.\ bit-flip channel as a closed-form polynomial in the flip rate \(p\); its coefficients are strictly determined by finite exhaustion of all \(64\) microstates:
\begin{itemize}
  \item \textbf{Generation script:} \texttt{\detokenize{scripts/exp_foldbin6_bitflip_stability_polynomial.py}}.
  \item \textbf{LaTeX generated output:} \texttt{\detokenize{sections/generated/eq_foldbin6_bitflip_stability_polynomial.tex}}.
  \item \textbf{Audit JSON:} \texttt{\detokenize{artifacts/export/foldbin6_bitflip_stability_polynomial.json}} (containing the stable count stratified by noise weight and the polynomial coefficients).
\end{itemize}

\paragraph{Verifiable interface: invertibility of the fiber-weighted coupling matrix \(M\) and the prime factorization of \texorpdfstring{$\det(M)$}{det(M)}}
Terminal proposition Theorem \ref{thm:terminal-window6-fiber-edge-coupling-det} constructs on \(X_6\) the weighted coupling matrix \(M\in\mathbb{Z}^{21\times 21}\) (inter-fiber edge count) induced by \(Q_6\) edges, and provides the row-sum observability identity and the integer factor decomposition certificate for \(\det(M)\neq 0\):
\begin{itemize}
  \item \textbf{Generation script:} \texttt{\detokenize{scripts/exp_window6_fiber_edge_coupling_matrix.py}}.
  \item \textbf{LaTeX generated output:} \texttt{\detokenize{sections/generated/eq_window6_fiber_edge_coupling_det.tex}}.
  \item \textbf{Audit JSON:} \texttt{\detokenize{artifacts/export/window6_fiber_edge_coupling_matrix.json}} (containing the full matrix \(M\), fiber sizes, and row-sum verification).
\end{itemize}

\paragraph{Verifiable interface: the \(1+8+12\) dictionary of \texorpdfstring{$X_6$}{X6} and the \(12=8+4\) two-tier heavy spectrum}
The ``\(1+8+12\)'' and ``\(12=8+4\)'' partition conclusions require no separate script:
they can be read directly from the three-row word table in Table \ref{tab:fold6_bin_uplift_choice_collapse}.
The first row gives the \(9\) words of \(S_{\mathrm{low}}\) (of which \(\texttt{000000}\) is the unique zero word and the remaining \(8\) form the \(1+8\) light block),
while the second and third rows combine to form the \(12\) heavy block, automatically split into two tiers by \(w_6=1\) (\(\Xi=2\), \(8\) words) and \(w_6=0,\ V_6(w)>8\) (\(\Xi=3\), \(4\) words).
From this one can directly compute the entropy weight \(h(w)=\log\Xi(w)\) and use it as partition weights for spectral gap alignment (\S\ref{subsec:group-unification-spectral-alignment}).
The corresponding verifiable outputs and scripts follow the generation chain for \(C_6\) orbits and dyadic uplift certificates described above (Table \ref{tab:fold6_bin_uplift_choice_collapse} and
\texttt{\detokenize{artifacts/export/window6_c6_orbit_patisalam_seed.json}}).

\paragraph{Verifiable interface: dual-scale selection law of the bin-fold degeneracy spectrum and the \texorpdfstring{$\dim=55$}{dim=55} unification anchor}
The dual-scale selection law of the bin-fold degeneracy spectrum uses the degeneracy spectrum
\(
d^{\mathrm{bin}}_m(w)=|\mathrm{Fold}^{\mathrm{bin}}_m{}^{-1}(w)|
\)
of the binary interval folding \(\mathrm{Fold}^{\mathrm{bin}}_m\)
as a third audit channel independent of the Zeckendorf--GUT signatures and the NAP selector:
at \(m=6\), the low-degeneracy sector reads out the counting interface \(12\), and the ghost capacity \(H^{\mathrm{bin}}_6=2^6-|X_6|=43\) completes to the unification dimension anchor \(55\);
at higher resolution \(m=10\), the minimal degeneracy sector count again exactly equals \(55\), forming a two-route locking.
The generation chain has been scripted and included in the one-click pipeline:
\begin{itemize}
  \item \textbf{Generation script:} \texttt{\detokenize{scripts/exp_fold_bin_degeneracy_spectrum_m6_m12.py}}.
  \item \textbf{LaTeX generated output:} \texttt{\detokenize{sections/generated/tab_fold_bin_degeneracy_spectrum_m6_m12.tex}} (Table \ref{tab:fold_bin_degeneracy_spectrum_m6_m12}).
  \item \textbf{Audit JSON:} \texttt{\detokenize{artifacts/export/fold_bin_degeneracy_spectrum_m6_m12.json}} (containing histograms for each \(m\), \(d_{\min}\), minimal sector size, and wall clock time).
\end{itemize}

\paragraph{Cross-audit interface: the cross-scale complement identity of \texorpdfstring{$\mathrm{Fold}^{\mathrm{bin}}$}{Foldbin}--\texorpdfstring{$\mathrm{Fold}$}{Fold}}
The interfaced conclusion \ref{par:gut-foldbin-fold-complement-identity} provides a cross-audit identity that directly locks ``the complement capacity of the most-rigid sector at low resolution'' with ``the extreme fiber modular scale at high resolution'';
the verification equations are independent derived outputs (not re-enumerating the spectral data) and have been included in the one-click pipeline:
\begin{itemize}
  \item \textbf{Generation script:} \texttt{\detokenize{scripts/exp_foldbin_fold_complement_identity_m6_m12.py}}.
  \item \textbf{LaTeX generated output:} \texttt{\detokenize{sections/generated/eq_foldbin_fold_complement_identity_m6_m12.tex}}.
  \item \textbf{Audit JSON:} \texttt{\detokenize{artifacts/export/foldbin_fold_complement_identity_m6_m12.json}} (containing per-\(m\) verification items and pass/fail flags).
\end{itemize}

\paragraph{Verifiable interface: four-step shift uplift isomorphism and the \(6+6+6+3\) root-type partition at \texorpdfstring{$m=10$}{m=10}}
The four-step shift uplift isomorphism canonically identifies \(X_{10}^{\mathrm{bdry}}\) with \(X_6\), and refines it via the \(B_3\) root dictionary into
\(6\) short roots, \(6\) \(A_2(+)\) roots, \(6\) \(A_2(-)\) roots, and \(3\) Cartan directions.
The corresponding \(21\)-word uplift dictionary and partition list are automatically generated by script:
\begin{itemize}
  \item \textbf{Generation script:} \texttt{\detokenize{scripts/exp_boundary_uplift_m10_b3c3_dictionary.py}}.
  \item \textbf{LaTeX generated outputs:} \texttt{\detokenize{sections/generated/eq_boundary_uplift_m10_isomorphism.tex}},
  \texttt{\detokenize{sections/generated/tab_boundary_uplift_m10_dictionary.tex}}.
  \item \textbf{Audit JSON:} \texttt{\detokenize{artifacts/export/boundary_uplift_m10_b3c3_dictionary.json}} (containing the finite-window bijection verification of \(\iota_n\), the \(m=10\) 21-word dictionary, and the \(A_2(\pm)\) partition list).
\end{itemize}

\begin{remark}[Double \texorpdfstring{$A_2$}{A2} sub-system splitting of long roots and sign rigidity]\label{rem:boundary-m10-longroot-double-a2-sign-rigidity}
Under the framework of Theorem \ref{thm:boundary-m10-b3c3-rigid-four-block-split},
the \(18\) root directions at \(m=10\) decompose strictly as
\(
6(\mathrm{short})\ \sqcup\ 6(A_2(+))\ \sqcup\ 6(A_2(-))
\), and the two sets \(A_2(\pm)\) are distinguished solely by the sign condition ``two nonzero components have the same/opposite sign.''
Therefore this splitting introduces no additional choice parameters at the dictionary level and can serve as an intrinsic bisection mechanism for \(\mathfrak{su}(3)\)-type sub-data.
\end{remark}

\paragraph{Verifiable interface: the NAP selector and the \texorpdfstring{$\ZZ_{34}$}{Z34} Fourier decomposition certificate at \texorpdfstring{$m=11$}{m=11}}
Subsection \ref{subsec:bdry-tower-zeck-gut} introduces, on top of the Zeckendorf--GUT signature dictionary, a new discrete selection law:
requiring that the non-abelian resonance layers \(\{6,8\}\) must be preserved as independent layers (NAP) and minimizing \(D=\dim\mathfrak g\) among all simple Lie algebras.
This selector discretely picks out \(SO(10)\) as the minimal solution and gives the ``\(4\to 11\)'' single signature substitution with the rigid counting interface \(34-1=33\).
Simultaneously, the \(34\) boundary modes at \(m=11\) carry a \(\ZZ_{34}\) translation action under the canonical numbering, thereby forcing a
\(34=1\oplus 33\) with \(33=1\oplus 16\times 2\) real representation decomposition template (providing auditable targets for subsequent projection gate/kernel spectral merging).
\begin{itemize}
  \item \textbf{NAP selector script:} \texttt{\detokenize{scripts/exp_nap_selector_so10.py}}.
  \item \textbf{NAP LaTeX generated outputs:} \texttt{\detokenize{sections/generated/tab_nap_selector_so10.tex}},
  \texttt{\detokenize{sections/generated/eq_nap_selector_so10.tex}}.
  \item \textbf{NAP audit JSON:} \texttt{\detokenize{artifacts/export/nap_selector_so10.json}} (containing scan parameters, candidate table, and minimal solution certificate).
  \item \textbf{\(\ZZ_{34}\) Fourier decomposition script:} \texttt{\detokenize{scripts/exp_m11_boundary_z34_fourier_decomposition.py}}.
  \item \textbf{\(\ZZ_{34}\) LaTeX generated output:} \texttt{\detokenize{sections/generated/eq_m11_boundary_z34_decomposition.tex}}.
  \item \textbf{\(\ZZ_{34}\) audit JSON:} \texttt{\detokenize{artifacts/export/m11_boundary_z34_fourier_decomposition.json}} (containing the canonical numbering \(j\mapsto (v_7,u_{11})\) and orthogonality/rotation check residuals).
\end{itemize}

\begin{remark}[Uniqueness of local uplift and the non-locality of \texorpdfstring{$SU(5)$}{SU(5)}/\texorpdfstring{$E_6$}{E6} branches]\label{rem:window6-local-uplift-excludes-su5-e6}
By Corollary \ref{cor:terminal-window6-local-uplift-admissibility},
under the local uplift that preserves the stable type label \(F_6\) and is realized by \(\mathrm{Aut}(Q_6)\), the only realizable integer displacement besides \(0\) is \(\pm 34\);
hence among the three branches \(\{F_8,F_9,F_{10}\}=\{21,34,55\}\), only \(F_9\) is compatible with this locality.
On the other hand, Theorem \ref{thm:terminal-window6-tail-three-branch} gives the top Zeckendorf terms of \((SU(5),SO(10),E_6)\) as aligned with \((F_8,F_9,F_{10})\) respectively.
Therefore, when the additional protocol constraint ``uplift must be realized by the local geometric stabilizer'' is imposed,
the top-level branches of \(SU(5)\) and \(E_6\) can only appear through non-local mechanisms, and the \(SO(10)\) branch is the unique candidate retaining local realizability.
\end{remark}

\begin{corollary}[``Why the Standard Model and not circular reasoning'': threefold rigidity closure (window-\(6\))]\label{cor:window6-threefold-rigidity-closure-sm-so10}
At the window-\(6\) anchor, simultaneously juxtapose the following three auditable inputs:
\begin{enumerate}
  \item \textbf{Geometric rigidity:} the local uplift admits only \(\delta=34\) besides \(0\) (Corollary \ref{cor:terminal-window6-local-uplift-admissibility}).
  \item \textbf{Boundary tower observability:} the minimal three-layer boundary tower gives the \((1,3,8)\) dimensional interface (\texttt{\detokenize{sections/generated/eq_boundary_tower_sm12.tex}}), rigidly selecting \(\mathfrak u(1)\oplus\mathfrak{su}(2)\oplus\mathfrak{su}(3)\) (Corollary \ref{cor:bdry-tower-minimal-resonance-triple-sm}).
  \item \textbf{Combinatorial rigidity:} \(\mathrm{Aut}(\Gamma_6)=\{1\}\) rules out nontrivial relabeling (Proposition \ref{prop:terminal-gamma6-rigidity} and Corollary \ref{cor:terminal-gamma6-multiplicity-rigidity}).
\end{enumerate}
Then the low-resolution visible gauge algebra is canonically localized up to isomorphism as
\[
\mathfrak g_{\mathrm{vis}}\ \cong\ \mathfrak u(1)\ \oplus\ \mathfrak{su}(2)\ \oplus\ \mathfrak{su}(3),
\]
without depending on any ``manually specifying the gauge group'' input.
Furthermore, when imposing the NAP constraint \(\{6,8\}\subseteq \Sigma(\mathfrak g)\) and minimizing the unification cost,
the scripted scan discretely locks the minimal simple completion to \(SO(10)\) and gives the single-layer substitution certificate
\(
\Sigma_{SO(10)}=\Sigma_{\mathrm{SM}}\setminus\{4\}\cup\{11\}
\)
with \(34-1=45-12\) (see \eqref{eq_nap_selector_so10});
the top-level term \(F_9=34\) of this substitution coincides with the unique integer fingerprint of the local uplift, thereby closing the non-circular inference chain of ``geometric locality--boundary tower minimality--unification completion'' on the same integer certificate.
\end{corollary}
\begin{proof}
The first part follows from juxtaposing Corollary \ref{cor:terminal-window6-local-uplift-admissibility}, Corollary \ref{cor:bdry-tower-minimal-resonance-triple-sm}, and Proposition \ref{prop:terminal-gamma6-rigidity} (with its corollaries).
The second part follows from the NAP selector generation equation \eqref{eq_nap_selector_so10}. \qed
\end{proof}

\paragraph{Verifiable check: the \texorpdfstring{$\mathrm{Sp}(6)$}{Sp(6)} isotropy point scan for the three-channel potential--response symplectic phase space}
Remark \ref{rem:sp6-minimal-continuous-envelope} provides a minimal framework for ``promoting the \(B_3/C_3\) discrete unification seed to a continuous envelope'': under the Legendre dual coordinates of the pressure \(P(\theta)\),
\((\theta;J=\nabla P)\) induces a standard symplectic form, and the minimal continuous group preserving the three-channel conjugate structure is \(\mathrm{Sp}(6)\).
A strictly verifiable corresponding task is: at a certain ``unification point,'' the second-order fluctuation scale should tend toward single-scale (isotropy), with the fully auditable figure of merit
\(
\mathcal{I}(\theta)=\log(\lambda_{\max}/\lambda_{\min})
\)
(see Appendix Remark \ref{rem:sync-kernel-3d-sp6-isotropy-scan}).
This scan has been scripted and included in the one-click pipeline:
\begin{itemize}
  \item \textbf{Generation script:} \texttt{\detokenize{scripts/exp_sync_kernel_weighted_sp6_isotropy_scan.py}}.
  \item \textbf{LaTeX generated output:} \texttt{\detokenize{sections/generated/tab_sync_kernel_weighted_sp6_isotropy_scan.tex}} (Table \ref{tab:sync_kernel_weighted_sp6_isotropy_scan}).
  \item \textbf{Audit JSON:} \texttt{\detokenize{artifacts/export/sync_kernel_weighted_sp6_isotropy_scan.json}} (containing the scan grid, \(\Sigma=\nabla^2P\), eigenvalues, and the optimal candidate point list).
  \item \textbf{Root-type split ratio script:} \texttt{\detokenize{scripts/exp_sp6_isotropy_root_split_ratio.py}} (reads the above JSON and computes \(R_{ls},R_{pm}\)).
  \item \textbf{Root-type split LaTeX generated output:} \texttt{\detokenize{sections/generated/tab_sp6_isotropy_root_split_ratio.tex}}.
  \item \textbf{Root-type split audit JSON:} \texttt{\detokenize{artifacts/export/sp6_isotropy_root_split_ratio.json}}.
\end{itemize}

\paragraph{Interface extension: reducing external time rescaling, folding defects, and coupling drift to the same primary scalar variable}
This chapter has already written the external time rescaling in Remark \ref{rem:lapse-modulated-tau-mix} as a minimal external time rescaling proxy
\(
N:=\kappa_0/\kappa
\)
(where \(\kappa\) is the background-dependent routing/inverse search overhead, in units of ``external cost per internal step'').
On the other hand, the folding side has already defined the fiber defect scalar
\begin{equation}\label{eq:gut_kappa_m_defect}
\kappa_m(\pi):=\mathbb{E}_{\pi}[\log d_m(X)]
\qquad(\text{Definition \ref{def:pom-kappa}; where }d_m(x)=|\Fold_m^{-1}(x)|).
\end{equation}
Therefore one can introduce a unified primary scalar variable (a pure definition, introducing no additional referents):
\[
u:=\log\frac{\kappa}{\kappa_0}=-\log N.
\]
The \emph{interface-level} assumption is: under spatially indexed readout (e.g., Hilbert-addressed windows), the local overhead field \(u(x)\) is isomorphic to some local folding statistic defect quantity,
for example reading out the local distribution \(\pi_{m,x}\in\Delta(X_m)\) as
\(
\kappa_m(\pi_{m,x})
\)
and setting
\[
u(x)\ \approx\ c\,\kappa_m(\pi_{m,x})
\qquad(c>0\ \text{is a calibration constant}).
\]
Once any effective coupling or effective constant drift is written as a function of the primary scalar (linearized form)
\[
g_i^{-2}(x)=a_i+b_i\,u(x)+\cdots,
\]
one obtains a \textbf{falsifiable correlation}:
\[
\Delta g_i^{-2}\ \approx\ b_i\,\Delta u\ =\ -b_i\,\Delta\log N.
\]
This correlation does not depend on a continuum limit; it only requires simultaneously outputting proxy quantities for \(N\) and \(g_i\) on the same data/same protocol family and performing a joint test.
\par
The more complete overhead\(\to\)external time rescaling\(\to\)linearization closure, as well as the executable protocol for reconstructing the spatial cost scalar variable from addressed windows and folding degeneracy statistics, belong to further constructions at the interface level;
this manuscript records here only its minimal isomorphism relation with the variables of this chapter and the falsifiable correlation form.

\paragraph{Finite verification template: Hecke rigidity and the finite prime skeleton decision via Sturm bounds}
\begin{theorem}[Finite decidability of auditable Langlands-type identification: Sturm boundedness \(\times\) finite-layer closure \(\times\) counterterm strictification]\label{thm:audit-finite-decidability-sturm-finite-layer}
Suppose that the readout of a candidate ``unified object'' is organized as a set of comparable coefficient/local factor data (e.g., the \(q\)-expansion coefficients of a modular form, Hecke eigenvalues \(a_p\), or the equivalent prime local factors/characteristic polynomials), and that its realization chain allows simultaneously outputting Dirichlet twist channels and finite-part drift budgets.
Then under the auditable framework of this paper, Langlands-type ``identification/matching'' (whether it belongs to a given Hecke eigenpacket/equivalence class) can be reduced to a \textbf{finite} verification problem, whose finiteness source can be explicitly decomposed into three components:
\begin{enumerate}
  \item \textbf{(Modular form side: Sturm boundedness)} Given weight \(k\) and level \(N\), the Sturm bound \(B_{k,N}\) gives a finite truncation decision: one need only verify up to \(B_{k,N}\) Fourier/Hecke data to decide equality (thereby replacing ``infinitely many conditions'' by a check on a finite prime skeleton).
  \item \textbf{(Dirichlet side: finite-layer closure)} Any finite projection word/finite realization's dependence on the Dirichlet axis bundle must close at some finite layer \(m_\ast\); hence ``full character scan/worst channel'' can be reduced to exhaustive verification over a finite set \(\widehat{G_{m_\ast}}\) (Proposition \ref{prop:pom-finite-modulus-stability}, see also Corollary \ref{cor:pom-local-finite-dirichlet-action}).
  \item \textbf{(Interface strictification side: counterterm cancellation of finite-part drift)} The commutation failure under the central extension framework falls entirely on the finite part/constant layer and can be strictly canceled by a purely constant counterterm gate, thereby promoting the residual commutation to a strict compatibility condition (Proposition \ref{prop:pom-counterterm-anom-cancel}, see also Theorem \ref{thm:pom-strictification-splitting}).
\end{enumerate}
Therefore, under the given error budget structure (finite sample/truncation/strictification), the above identification problem is no longer an open-ended infinite verification but a finite decision problem whose scale is jointly controlled by \((B_{k,N},m_\ast)\) and the finite-part counterterm budget.
\end{theorem}

\begin{remark}[Summary conclusion: the closed main line of lapse--Dirichlet--Adams--counterterm--Sturm]\label{rem:audit-summary-lapse-dirichlet-adams-sturm}
This manuscript systematically compresses ``auditable computation'' into finite certificate chains: lapse (measure transformation/external time rescaling) can be internalized as a weighted periodic readout and gives explicit finite sample upper bounds (Theorem \ref{thm:op-algebra-lapse-weighted-wish-koksma}); Dirichlet twist channels are locally finite under finite-layer closure and can be exhaustively verified (Proposition \ref{prop:pom-finite-modulus-stability} and Corollary \ref{cor:pom-local-finite-dirichlet-action}); the Adams--Dirichlet arithmetic scaling gate unifies the \((\chi,u)\) power semigroup action as a renormalization monad (Definition \ref{def:witt-adams-dirichlet-renorm-gate} and Proposition \ref{prop:witt-renorm-monoid}); anomalies are strictified to controlled constant-layer drift via central extension splitting and counterterm strictification (Theorem \ref{thm:pom-strictification-splitting}).
On this basis, Sturm boundedness grounds Hecke rigidity as a finite prime skeleton decision (Theorem \ref{thm:audit-finite-decidability-sturm-finite-layer}), thereby forming a closed structural main line that can be cited item by item, without manual numbering.
The existing script \texttt{\detokenize{scripts/exp_hecke_e4_prime_audit.py}} provides the minimal reproducible instance of this template: in the zero-parameter baseline case, using the Hecke coefficient table/JSON output at finitely many prime points as a rapid negation test.
\end{remark}

\begin{corollary}[Finite prime skeleton negation criterion under Sturm--Hecke boundedness: three primes suffice to exclude the scalar \texorpdfstring{$X(1)$}{X(1)} weight \texorpdfstring{$4$}{4} case]\label{cor:hecke-e4-three-prime-refutation}
If the top-level unified object is declared to be a \emph{level-1 (\(X(1)\)) scalar weight \(4\) Hecke-closed} modular form,
then this space is one-dimensional and spanned by the Eisenstein series \(E_4\), so its prime Hecke coefficients must satisfy
\[
\boxed{\ a_p=240(1+p^3)\qquad(p\ \text{prime})\ }.
\]
Thus under the minimal negation framework, one need only compare coefficients at \(p\in\{2,3,5\}\) to form a rapid negation test:
any significant deviation implies that ``the top-level unification cannot remain at the \(X(1)\) scalar one-dimensional + Hecke rigidity level and must enter the vector-valued/congruence level lifting branch.''
\end{corollary}
\par
We package the minimal negation criterion at the prime points \(p\in\{2,3,5\}\) as a numerical table and JSON for reproducible output (serving as a ``zero-fit necessary condition'' rapid negation test):
\begin{itemize}
  \item \textbf{Generation script:} \texttt{\detokenize{scripts/exp_hecke_e4_prime_audit.py}}.
  \item \textbf{LaTeX generated output:} \texttt{\detokenize{sections/generated/tab_hecke_e4_prime_audit_p235.tex}}.
  \item \textbf{Audit JSON:} \texttt{\detokenize{artifacts/export/hecke_e4_prime_audit_p235.json}}.
\end{itemize}
